\chapter{Forward entropy at the fine horizon}\label{u:sec:fine}

The goal of this section is to prove \cref{u:prop:fine-kl}, which bounds
the relative entropy of the fine proposal with respect to $\Pi_n$.
We use the source and replay laws of \cref{u:def:laws}, the scalar
identities of \cref{u:prop:scalar,u:lem:cross}, and the quadrature estimate
of \cref{u:lem:quad}. We use this bound at the two ends of the
connecting sequence in \cref{u:sec:connections}.
For the reference measure and the comparison of likelihoods, we follow
\cite[Section~4, Proposition~5]{oai-typical-upper}. We then estimate the
expected energy directly, including the outcomes on which the query uses
the deterministic fallback.

Throughout this section, $V=t_{N+1}$ is the terminal time of a proposal
with $N<n$ queries. We write $m_i(t)=m(t,Y_i(t))$ and
$\chi_i(t)=\chi(t,Y_i(t))$. Expectations with subscript $\Lambda$
refer to the linked source or replay law, according to the space of
the integrand; these laws agree on the complete row histories and on
the output by \cref{u:sec:proposal-laws}.

\section{The reference measure}
\label{u:sec:reference}

On the unlinked replay space, we regard $(W,C,B^0,\mathrm{br},x)$ as static.
The $D_i$ are independent Brownian motions, independent of these
static data. The replay field at time $t$ is a Borel function of these
static data and $D_{[0,t]}$: each next query reads the field only at
the left endpoint of its new cell. We define
\begin{equation}\label{u:eq:reference-coefficient}
 \xi_i(t)=\frac{x_i \chi(t,Y_i(t))}{1+x_i m(t,Y_i(t))},\qquad
 |\xi_i(t)|\le2.
\end{equation}
The coefficient $\sigma(t)\xi_i(t)$ is continuous and predictable in this replay filtration.
It vanishes before $v_*$; afterward it is $\sigma(t)(x_i-m_i(t))$.
Under $\mu^{\rm rep}_0$, let
\begin{equation}\label{u:eq:reference-density}
 \begin{aligned}
 \mathscr M_D&=\sum_i\int_0^V\sigma(t)\xi_i(t)\dd D_i(t),\\
 \langle\mathscr M_D\rangle
   &=\sum_i\int_0^V\sigma(t)^2\xi_i(t)^2\dd t\le4nV,\\
 E_D&=e^{\mathscr M_D-\langle\mathscr M_D\rangle/2},\qquad
 R_D=E_D\mu^{\rm rep}_0.
 \end{aligned}
\end{equation}
We define $\mathscr M_D$ as the $L^2$ limit of the vector of
predictable sums at the left endpoints of dyadic partitions. This
definition, and not the normalization alone, is what identifies it
with the sum of the row integrals below. We choose versions according to the
conventions of \cref{u:sec:setup}.

\trilogyneedspace{12\baselineskip}
\begin{lemma}[The reference measure and the row integrals]
\label{u:lem:reference}
The measure $R_D$ is a probability measure, preserves the complete
static marginal, and has output marginal $\gamma_n\otimes \Unif_n$ on $(g,x)$.
For one unlinked row, let $\xi(t)=x\chi_t/(1+xm_t)$ and
\[
 \begin{aligned}
 \mathscr M_D^{\rm row}
   &=\int_0^V\sigma(t)\xi(t)\dd D(t),\\
 \langle\mathscr M_D^{\rm row}\rangle
   &=\int_0^V\sigma(t)^2\xi(t)^2\dd t.
 \end{aligned}
\]
The integral denotes its unlinked $L^2$ left-sum version;
we use the same Borel version under the equivalent linked row law.
Moreover,
\begin{equation}\label{u:eq:reference-ratio}
 \begin{aligned}
 \frac{\dd\Lambda^{\rm rep}}{\dd R_D}&=L/E_D,\qquad
 I^{(n)}=\log(L/E_D)=\sum_{i=1}^n I_{V,i},\\
 I_V&=\log(1+x m_V)-\mathscr M_D^{\rm row}
                   +\tfrac12\langle\mathscr M_D^{\rm row}\rangle,
 \end{aligned}
\end{equation}
Here $I_{V,i}$ is $I_V$ evaluated on row $i$.
The row histories are independent with law $\nu_V$ of
\cref{u:sec:scalar}, so the summands are independent with
common law $(I_V)_*\nu_V$.
Finally,
\begin{equation}\label{u:eq:reference-mean}
 \E_{\Lambda^{\rm rep}}I^{(n)}=n j(V).
\end{equation}
\end{lemma}
\begin{proof}
We first normalize the exponential conditional on the static data.
We then transport its defining left sums to the source row space and
compute the mean in the source natural filtration.

\paragraph{Normalization.}
This is the conditional stochastic-exponential normalization used in
\cite[Section~4, proof of Proposition~5]{oai-typical-upper}; the
bounded-integrand change of measure is Girsanov's theorem
\cite[Theorem~3.5.1]{karatzas-shreve-1991}.
Fix the static data and a deterministic grid for the left sums. The last
increment of $D$ is Gaussian and independent of the earlier ones, and
integrating it out shows that its factor in the exponential, together with
the compensator, has conditional mean one. Iterating over the increments, we find
that the exponential of the left sums has mean one. Since the predictable
coefficients are bounded and continuous, their left approximants
converge in $L^2(\dd t\,\dd\mu^{\rm rep}_0)$. By isometry, the sums
then converge to the integral defining $\mathscr M_D$. The exponentials are uniformly
integrable, because the same Gaussian identity at twice the coefficient
bounds their second moments by $e^{4nV}$. Testing against a bounded
function of the complete static data before taking the limit, we
conclude that the limiting exponential has conditional mean one given
these data.
This proves that $R_D$ is normalized and that it preserves the
complete static marginal, including the output marginal.

With the Borel versions fixed above, and after a modification on a null set,
$E_D$ is positive and finite everywhere.
Since $L>0$, multiplying the densities gives
$R_D\cdot(L/E_D)=\mu^{\rm rep}_0\cdot L=\Lambda^{\rm rep}$.

\paragraph{Identification of the row integrals.}
To identify the integral, we use the diagram of laws and the history
identity in \cref{u:sec:proposal-laws} for the unlinked laws. At every
finite dyadic level, the pullback of the vector sum on the replay
space equals the sum of the corresponding functions of the row
histories $(x,Y,B,D)$. On the unlinked
row space, $D$ is a Brownian motion in the filtration generated by the two original
drivers and by the independent uniform sign, which is revealed at time zero. The
coefficient $\sigma\xi_i$ from \eqref{u:eq:reference-coefficient} is bounded and
predictable, so the same sums converge in $L^2$ to its integral. The law identity and the
Cauchy--Schwarz inequality for finite sums imply that the sum of the row limits
converges in $L^2$ on the replay space. By uniqueness of the limit, $\mathscr M_D$ is
the sum of the row integrals. By Fubini's theorem in the time variable,
the brackets split in the same way.
This proves \eqref{u:eq:reference-ratio} almost everywhere under the
unlinked law. The linked and unlinked laws are equivalent, which implies that the
identity also holds almost everywhere under the linked law. Here the transported variables are the complete histories, and the
versions are fixed at the given $n,V$.

\paragraph{The mean.}
To compute the mean, we work in the natural filtration of the unlinked source,
$\mathcal F_t=\sigma(x,(B_i,D_i)_{[0,t]}:i\le n)$, which does not contain $W$,
$d_j$, and $W_0$. By \cref{u:lem:cross}, we have
$\dd m_i=\chi_i(\rho\dd B_i+\sigma\dd D_i)$. Using It\^o's formula,
localized where $1+x_i m_i\ge1/l$ and then letting $l\to\infty$, we find
\[
 \log L=\sum_i\int_0^V\xi_i(\rho\dd B_i+\sigma\dd D_i)
                         -\tfrac12\sum_i\int_0^V\xi_i^2\dd t.
\]
The localization can be removed because $|\xi_i|\le2$ and the denominator is
strictly positive on every finite continuous sample path.
The integral against $D$ is the pullback identified above, so
\begin{equation}\label{u:eq:information-brownian}
 I^{(n)}=\sum_i\int_0^V\rho\xi_i\dd B_i
                        -\tfrac12\sum_i\int_0^V\rho^2\xi_i^2\dd t.
\end{equation}
The process $L_t=\prod_i(1+x_i m_i(t))$ is the positive bounded
martingale that defines the terminal link, with $L_0=1$.
The expectation in Novikov's condition is at most $e^{2nV}$.
By Girsanov's theorem for bounded integrands
\cite[Theorem~3.5.1]{karatzas-shreve-1991}, under $\Lambda^{\rm src}$
the processes
\[
 B_i^\Lambda=B_i-\int_0^\cdot\rho\xi_i\dd t,\qquad
 D_i^\Lambda=D_i-\int_0^\cdot\sigma\xi_i\dd t
\]
form a Brownian vector in this natural filtration.
It follows from \eqref{u:eq:information-brownian} that
\[
 \E_\Lambda I^{(n)}
   =\frac12\sum_i\E_\Lambda\int_0^V\rho^2\xi_i^2\dd t.
\]
Given the past of the spin-free pair at time $t$, the linked spin probabilities
are $(1+x m_t)/2$, and therefore
\[
 \sum_{x=\pm1}\frac{1+x m_t}{2}
       \frac{\chi_t^2}{(1+x m_t)^2}=\frac{\chi_t^2}{1-m_t^2}.
\]
The spin-free paths keep their original law. Since the integrands
are bounded, we may exchange the expectation and the time integral. Using the definition of $j$ in \cref{u:prop:scalar}, we obtain
\eqref{u:eq:reference-mean}.

\paragraph{Integrability.}
All signed quantities used here are integrable under the indicated
laws. By isometry, we have $\mathscr M_D\in L^2(\mu^{\rm rep}_0)$; also
$\langle\mathscr M_D\rangle\le4nV$.
For fixed $n$, the bound $L\le2^n$ transfers this $L^2$ integrability to
$\Lambda$. Also, $\log L\in L^2(\Lambda)$, since $t(\log t)^2$ is bounded
on $0<t\le2$ and all other spin factors average to one.
After compensation in \eqref{u:eq:information-brownian}, we obtain the bound
$\E_\Lambda(I^{(n)})^2\le8nV+8n^2V^2$.
The calculation with conditional exponentials above also proves the case of
Girsanov's theorem with bounded integrands, by exponential tilting of
finitely many Gaussian steps and passage to the limit in $L^2$. Throughout,
the filtration does not contain the final matrix.
\end{proof}

\begin{lemma}[Forward relative entropy]\label{u:lem:forward-kl}
For the output $Q_V$ of \cref{u:def:laws},
\begin{equation}\label{u:eq:forward-kl}
 \KL(Q_V\Vert \Pi_n)\le n j(V)-n\log2+\E_{\gamma_n}F^c
                               -\beta\E_\Lambda H^c_g(x).
\end{equation}
\end{lemma}
\begin{proof}
We use the lifted-reference and entropy-contraction argument of
\cite[Section~4, proof of Proposition~5]{oai-typical-upper}.
On the replay space, we define
$\dd\widehat P=2^n e^{\beta H^c_g(x)-F^c(g)}\dd R_D$.
Its total mass is one, since $R_D$ projects to $\gamma_n\otimes\Unif_n$;
its output law is $\Pi_n$. Further,
\[
 \log\frac{\dd\Lambda}{\dd\widehat P}
      = I^{(n)}-n\log2+F^c(g)-\beta H^c_g(x).
\]
Averaging pointwise over the spins, we see that the matrix marginal
of the linked law is the GOE law. In particular, with $H_*(W)=\max_x|x^{\mathsf T}Wx/2|$,
\begin{equation}\label{u:eq:energy-integrability}
 \E_\Lambda H_*^2\le\frac{n^2}{4}\E_{\mathcal G_n}\|W\|_F^2
                =\frac{n^2(n+1)}4.
\end{equation}
Also $|F^c|\le n\log2+|\beta|H_*$. These bounds and \cref{u:lem:reference}
justify the finite KL calculation. The data processing inequality for the output
map \cite[Chapter~4, Exercise~4.10]{boucheron-lugosi-massart-2013}
and \eqref{u:eq:reference-mean} imply \eqref{u:eq:forward-kl}.
\end{proof}

\section{Fine energy}
\label{u:sec:fine-energy}

We now estimate the expected energy, one query at a time.
The decomposition into queried columns, diagonal innovations and an
unrevealed GOE block is the energy decomposition in
\cite[Section~4, proof of Proposition~5]{oai-typical-upper}. The
conditional expectation calculation below also controls the fallback
outcomes and the unequal cell durations.
Let $\mathcal F_{t_j}^{\rm paths}$ be the natural filtration of the
spin-free pair of drivers in the source, through time $t_j$.
Under the linked law, the marginal law of these paths is unchanged. The
conditional laws of the future rows and of the signs factor across
coordinates, because the unlinked rows factor, their pasts factor, and
the link is a product. We obtain their conditional densities by
integrating the terminal link of each row over its independent future.
In particular, the pairs $(e_{j,i},x_i)$
are conditionally independent across $i$ and satisfy
\begin{equation}\label{u:eq:fine-conditional-correlation}
 \begin{aligned}
 \E[e_{j,i}\mid\mathcal F_{t_j}^{\rm paths}]&=0,\qquad
 \E[x_i\mid\mathcal F_{t_j}^{\rm paths}]=m_i(t_j),\\
 \psi_i=\E[e_{j,i}x_i\mid\mathcal F_{t_j}^{\rm paths}]
   &=\frac1{\sqrt{h_j}}\int_{t_j}^{t_{j+1}}\rho(v)
            \E[\chi_i(v)\mid\mathcal F_{t_j}^{\rm paths}]\dd v.
 \end{aligned}
\end{equation}
The last identity is the cross-moment identity of \cref{u:sec:scalar}
tested against functions of the past. Here we condition on the past of
the original drivers, not on $W$. The Picard map with fresh future
drivers gives a Borel version of $\psi_i$ that depends only on the past of
its own row, in accordance with the conventions on versions in \cref{u:sec:setup}.
If $h_{j-1},h_j\le1/8$, then the relative variation
estimates of \cref{u:prop:scalar} and the growth of the schedule imply that
\begin{equation}\label{u:eq:fine-correlation-bound}
 0\le\psi_i\le\beta\sqrt{h_jr(t_j)}\le2\beta s_j,
 \qquad \E\overline\psi=b_j,\quad
 \overline\psi=n^{-1}\sum_i\psi_i.
\end{equation}
The factor two accounts for the different durations of the previous
and current cells.

In this section, we write $\varrho_j=\normn{P_{j-1}\mathcal T_j}$ for the
residual norm. Whether the query is normalized or uses the fallback, we have
$\ipn{u_j}{m(t_j)}=s_j\varrho_j$, because
$m(t_j)=\sum_{l\le j}s_l\mathcal T_l$ and each $\mathcal T_l$ belongs to the span
of the first $l$ queries. Also $P_jm(t_j)=0$.

\trilogyneedspace{12\baselineskip}
\begin{lemma}[Exact one-query conditional identity]\label{u:lem:onequery}
With $c_j=\ipn{u_j}x$,
\begin{equation}\label{u:eq:onequery}
 \begin{aligned}
 &\E[c_j\ipn{x}{P_je_j}\mid\mathcal F_{t_j}^{\rm paths}]\\
 &\qquad=s_j\varrho_j\frac1n\sum_i(P_j)_{ii}\psi_i
       -\frac2{n^2}\sum_i u_{j,i}(P_j)_{ii}\psi_i m_i(t_j).
 \end{aligned}
\end{equation}
\end{lemma}
\begin{proof}
We expand $(n^{-1}\sum_i u_i x_i)(n^{-1}\sum_{k,l}x_k P_{kl}e_l)$.
Since the row pairs are conditionally independent and $x_i^2=1$, the mixed
third moment is
\[
 \E[x_i x_k e_l\mid\mathrm{past}]
 =\mathbf1_{l=i}\psi_i m_k+\mathbf1_{l=k}m_i\psi_k
                  -2\mathbf1_{i=k=l}m_i\psi_i.
\]
For $i=k$ it is zero by Gaussian centering. For $i\ne k$, it follows
from independence, also when $l$ differs from both indices. The
first contraction contains $P_jm(t_j)=0$, and the second contains
$\ipn{u_j}{m(t_j)}=s_j\varrho_j$.
In more detail, given the past of a row, a fresh seed for the row consists
of the current $B$ bridge, the future $B$ path after that cell,
the future $D$ increments, and a terminal uniform threshold.
It is independent of the standard Gaussian endpoint $e_{j,i}$;
these seeds and endpoints are independent across coordinates.
Since the Picard solution is causal, the terminal sign is a Borel
function of its own endpoint and seed. These conditional product laws justify the
third-moment calculation above.
All terms are integrable, since the signs are bounded and Gaussian
first moments are finite.
\end{proof}

\begin{lemma}[Expected energy at the fine horizon]\label{u:lem:fine-energy}
For $V=t_{N+1}$, $\max h_j\le1/8$,
\begin{equation}\label{u:eq:fine-energy}
 \E_\Lambda\frac{H^c_g(x)}n\ge e(V)-C_\beta\left[
 h+\frac{(N+1)(1+V)+\sqrt N}{n}
                         +\frac{(N+1)^2VR_V}{n^2}\right].
\end{equation}
\end{lemma}
\begin{proof}
We first estimate the expectation in \cref{u:lem:onequery}, including
the residual norm and the projection diagonals. We then sum over the
queries and use the quadrature estimate of \cref{u:lem:quad}.

\paragraph{The residual norm.}
Since $\sum_i|u_{j,i}|\le n$, the correction in \eqref{u:eq:onequery} is at
most $4\beta s_j/n$ in absolute value. Removing the projection diagonals
from its first term costs at most $2\beta s_j^2j \varrho_j/n$, and
$\E \varrho_j\le\sqrt{\E\normn{\mathcal T_j}^2}=1$.
For this calculation, we write
$Q=\normn{\mathcal T_j}^2$, $Y=\normn{U_{j-1}\mathcal T_j}^2$, $X=\varrho_j^2=Q-Y$.
For $X\ge0$,
\begin{equation}\label{u:eq:residual-square-root}
 \sqrt X\ge1+(X-1)/2-(X-1)^2/2.
\end{equation}
Indeed, the remainder in the tangent-line identity for the square root
is $(\sqrt X-1)^2/2\le(X-1)^2/2$.
Since the row pairs $(\psi_i,\mathcal T_{j,i})$ are independent and identically distributed, we have
\begin{equation}\label{u:eq:residual-covariance}
 \E[\overline\psi(Q-1)]
       =n^{-1}\operatorname{Cov}(\psi_i,\mathcal T_{j,i}^2)
       \ge-2\beta s_j/n.
\end{equation}
Here we used that $\E \mathcal T_{j,i}^2=1$ and $\psi_i\ge0$.

Let $h_-=h_{j-1}$. Conditional on the past of the original pair up to time
$t_{j-1}$, the coordinates of the new vector $\mathcal T_j$
are independent, centered martingale increments, with conditional
second moment at most $h_-/s_j^2$ and fourth moment at most
$3h_-^2/s_j^4$. The bound on the fourth moment follows from It\^o's formula, applied to
the fourth power of an increment with coefficient $\chi\le1$, and the
bound on the second moment. The rank-$(j-1)$ projection $U_{j-1}$
is measurable at this time. Expanding its quadratic form using
conditional centering and independence, we obtain
\begin{equation}\label{u:eq:residual-projection-moments}
 \E Y\le\frac{j h_-}{ns_j^2},\qquad
 \E Y^2\le\frac{3(j^2+j)h_-^2}{n^2s_j^4}.
\end{equation}
For the unconditional fourth moment, we use the elementary bound
$\E Z_T^4\le144\E\langle Z\rangle_T^2$ for a martingale starting at
zero. Indeed, by It\^o's formula, we have
\[
 \E Z_T^4\le6(\E\sup|Z|^4)^{1/2}
                 (\E\langle Z\rangle_T^2)^{1/2}.
\]
Doob's fourth-moment inequality then implies the stated bound with
constant $144$.
Since $\langle\Delta m\rangle\le h_-$ and its expectation is $s_j^2$, we have
\begin{equation}\label{u:eq:residual-fourth-moments}
 \E \mathcal T_{j,i}^4\le144h_-/s_j^2,\qquad
 \E(Q-1)^2\le144h_-/(ns_j^2).
\end{equation}
All conditional statements here refer only to the pasts of the original drivers.

\paragraph{Summation of the query energies.}
We multiply \eqref{u:eq:residual-square-root} by $s_j\overline\psi\ge0$ and use \eqref{u:eq:fine-correlation-bound}, \eqref{u:eq:residual-covariance}--\eqref{u:eq:residual-fourth-moments},
and $(X-1)^2\le2(Q-1)^2+2Y^2$. Inserting the result into \eqref{u:eq:onequery}, we obtain
\begin{equation}\label{u:eq:onequery-energy-bound}
 \begin{aligned}
 \E[c_j\ipn{x}{P_je_j}]\ge s_jb_j-C_\beta\biggl[
 &\frac{(j+1)s_j^2+s_j+(j+1)h_-}{n}\\
 &+\frac{(j^2+j)h_-^2}{n^2s_j^2}\biggr].
 \end{aligned}
\end{equation}
This derivation also covers a vanishing residual, $\varrho_j=0$.
When summing over the queries, we use $\sum s_j^2\le1$, $\sum s_j\le\sqrt N$,
$s_j^2\ge h_-/R_V$, and $\sum h_-\le V$.

Expanding the assembled matrix of \cref{u:sec:proposal-schedule},
we have, on every outcome of the source,
\begin{equation}\label{u:eq:fine-energy-decomposition}
 \begin{aligned}
 \frac{H^c_g(x)}n&=E_{\rm qry}+Z,\qquad
 E_{\rm qry}=\sum_j c_j\ipn{x}{P_je_j},\\
 Z&=\frac12\sum_j c_j^2d_j
       +\frac1{2n}(P_Nx)^{\mathsf T}W_0(P_Nx).
 \end{aligned}
\end{equation}
Under the linked law, the unused variables $d_j,W_0$ are independent of
the complete scalar paths and the spins, because the link does not depend on them.
Thus $\E[Z\mid\text{complete scalar paths and spins}]=0$.
The required integrability follows from the Gaussian moments of these variables and \eqref{u:eq:energy-integrability}.
Summing \eqref{u:eq:onequery-energy-bound} and using
\cref{u:lem:quad} proves \eqref{u:eq:fine-energy}.
\end{proof}

\trilogyneedspace{12\baselineskip}
\begin{proposition}[Forward entropy at the fine horizon]\label{u:prop:fine-kl}
Let $h=n^{-\ex a}$, $V=V_{\ex f}$ as in \cref{u:lem:quad}, with
\begin{equation}\label{u:eq:fine-parameters}
 \ex a>\ex f/4,\quad \ex a+\ex f/2<1,\quad0<\ex k_{\rm f}<\min(\ex a,\ex f,1-\ex a).
\end{equation}
For the primitive fine output $Q_F$,
\begin{equation}\label{u:eq:fine-kl}
 \KL(Q_F\Vert \Pi_n)\le C_\beta n^{1-\ex k_{\rm f}}.
\end{equation}
The disorder and spin marginals of $Q_F$ are $\gamma_n$ and
$\Unif_n$, respectively.
\end{proposition}
\begin{proof}
Combining the entropy--energy identity of \cref{u:prop:scalar}, proved in
\cite[Lemma~4]{oai-typical-upper}, with
\eqref{u:eq:forward-kl} and \eqref{u:eq:fine-energy}, we find
\begin{equation}\label{u:eq:fine-kl-budget}
 \begin{aligned}
 \KL(Q_F\Vert\Pi_n)\le C_\beta\biggl[
 &nh+ne^{-V}+(N+1)(1+V)+\sqrt N\\
 &+\frac{(N+1)^2VR_V}{n}\biggr].
 \end{aligned}
\end{equation}
Here we also used $\E_{\gamma_n}F^c\le nP_\beta$ and
$p(V)^2\le C_\beta e^{-V}$. By \cref{u:lem:quad,u:lem:tails},
the powers of the five terms, up to fixed logarithmic factors, are
$1-\ex a$, $1-\ex f$, $\ex a$, $\ex a/2$, and $2\ex a+\ex f/2-1$.
The inequality $\ex a+\ex f/2<1$ implies that the last power is
smaller than $\ex a$.
Since the inequalities in \eqref{u:eq:fine-parameters} are strict, they absorb the constants and logarithmic factors.
For the affine family of \cref{u:sec:budget}, these powers are $1/2,1/2,1/2,1/4,1/4$, and the claimed
bound is $C_\beta n^{1/2+2\varepsilon_\omega}$. The energy estimate holds without restriction to a Gram event,
including every outcome on which a query uses the fallback. The disorder
marginal follows from averaging over the spins. The spin
marginal is $\Unif_n$ because the linked scalar rows are independent
and, by symmetry, $\E m_V=0$ in each row.
\end{proof}

